\documentclass[11pt]{article}
\usepackage[utf8]{inputenc}
\usepackage[T1]{fontenc}
\usepackage{lmodern}
\usepackage{amsmath,amssymb,amsthm}
\usepackage{enumitem}
\usepackage[margin=1in]{geometry}
\usepackage[hidelinks]{hyperref}
\setlength{\emergencystretch}{3em}

\newtheorem{theorem}{Theorem}
\newtheorem{corollary}[theorem]{Corollary}
\theoremstyle{definition}
\newtheorem{remark}[theorem]{Remark}
\newcommand{\G}{\mathcal G_{\Sigma}}
\newcommand{\C}{\mathcal C_{\Sigma}}
\newcommand{\Sel}{\operatorname{Sel}_{\chi}}
\newcommand{\Nat}{\operatorname{Nat}}
\newcommand{\Aut}{\operatorname{Aut}}
\newcommand{\Img}{\operatorname{im}}

\title{EQSRC-NO-NATURAL-SELECTOR-NONUNIQUENESS-THEOREM-V1\\
General no-selector and nonuniqueness theorems on the fixed certified core}
\author{The Aether-Flow Research Project---draft/control}
\date{2026-07-21}

\begin{document}
\maketitle

\begin{abstract}
On the unchanged proposal-only certified-isomorphism core fixed by P2-T01,
this packet classifies deterministic source-natural selectors as a product of
automorphism fixed loci. It proves the empty-fixed obstruction, gives the
globally quantified nonuniqueness criterion, separates selector multiplicity
from multiplicity of induced relation sections, and states the extra equalizer
equations required by noninvertible arrows. The proof is complete relative
to the declared core, functor, component-coordinate data, and set-theoretic
assumptions. It does not derive or adopt those inputs from current ontology.
\end{abstract}

\section{Authority and fixed domain}

This artifact is \texttt{draft/control}, \texttt{proposal-only}, and
\texttt{source-extension data}. Its adoption status is
\texttt{blocked\_adoption\_open\_continuation}. The source category,
eligible-choice functor, certified core, and relation interface are exactly
the P2-T01 objects inspected by P2-T02; they are fixed before selector
evaluation. Proof coordinates introduced below are mathematical coordinates
for the classification and are not added source marks.

Current ontology does not derive or adopt \(\C\), \(\G\), \(S_\chi\),
\(K\), a selector, the component-coordinate data, or any selector-relevant
extension. Structural automorphisms in \(\G\) are not thereby physical
gauge transformations. Both scientific ledgers remain unchanged.

\section{Declared objects and assumptions}

Let \(\G=\operatorname{Core}(\C)\) be the small, wide groupoid of the
unchanged P2-T01 proposal category, and let
\[
  S_\chi:\G\longrightarrow \mathbf{Set}
\]
be its unchanged eligible-choice functor. Let \(I=\pi_0(\G)\) be a set of
connected components. Supply one representative \(r_i\) for each
\(i\in I\), and for every object \(X\) in component \(i=[X]\), supply a
certified isomorphism
\[
  \tau_X:r_i\longrightarrow X,
  \qquad \tau_{r_i}=1_{r_i}.
\]
Define
\[
  H_i:=\Aut_{\G}(r_i),
  \qquad
  F_i:=S_\chi(r_i)^{H_i}
      =\{x\in S_\chi(r_i):S_\chi(h)x=x\text{ for all }h\in H_i\},
\]
and
\[
  \Sel(\G):=\Nat(1_{\G},S_\chi).
\]
The complete machine-readable assumption ledger is stored beside this source.
The representative and transport data can be replaced by any other supplied
coordinates without changing the reconstructed selector: the difference is
an automorphism that fixes the representative value.

\section{Core classification, obstruction, and nonuniqueness}

\begin{theorem}[Complete certified-core selector classification]
\label{thm:classification}
Evaluation at the supplied component representatives is a bijection
\[
  \operatorname{ev}:\Sel(\G)\xrightarrow{\;\cong\;}
  \prod_{i\in I}F_i,
  \qquad
  \sigma\longmapsto (\sigma_{r_i})_{i\in I}.
\]
Its inverse sends \(x=(x_i)_{i\in I}\) to
\[
  \sigma^x_X:=S_\chi(\tau_X)(x_{[X]}).
\]
Consequently:
\begin{enumerate}[label=(\alph*)]
  \item \(\Sel(\G)=\varnothing\) if and only if
        \(\prod_{i\in I}F_i=\varnothing\).
  \item If any \(F_i\) is empty, no deterministic source-natural selector
        exists on \(\G\).
  \item If every \(F_i\) is nonempty, a selector exists exactly when a
        product element (equivalently, a choice function for this family) is
        available. Finite \(I\), a supplied choice tuple, or an explicitly
        assumed suitable axiom of choice is sufficient.
  \item \(\Sel(\G)\) has exactly one element if and only if every \(F_i\)
        is a singleton.
  \item \(\lvert\Sel(\G)\rvert\geq2\) if and only if there is an
        \(i\in I\) with \(\lvert F_i\rvert\geq2\) and
        \(\prod_{j\neq i}F_j\neq\varnothing\).
\end{enumerate}
\end{theorem}

\begin{proof}
If \(\sigma:1_{\G}\Rightarrow S_\chi\) is natural and \(h\in H_i\),
naturality at \(h:r_i\to r_i\) gives
\(S_\chi(h)(\sigma_{r_i})=\sigma_{r_i}\). Thus evaluation lands in the
stated product.

Conversely, let \(x\in\prod_iF_i\) and define \(\sigma^x\) by the displayed
formula. For any \(f:X\to Y\) in \(\G\), the two objects lie in one
component \(i\), and
\[
  a:=\tau_Y^{-1}f\tau_X\in H_i.
\]
Since \(x_i\in F_i\),
\[
  S_\chi(f)\sigma^x_X
  =S_\chi(f\tau_X)x_i
  =S_\chi(\tau_Y a)x_i
  =S_\chi(\tau_Y)x_i
  =\sigma^x_Y.
\]
Hence \(\sigma^x\) is natural. At \(r_i\), the identity transport gives
\(\sigma^x_{r_i}=x_i\). In the other direction, naturality of \(\sigma\)
at \(\tau_X\) gives
\(\sigma_X=S_\chi(\tau_X)(\sigma_{r_i})\). The two maps are inverse.

Parts (a)--(c) follow from the bijection and the definition of a Cartesian
product. This wording is deliberate: in ZF, arbitrary factorwise
nonemptiness need not supply a product element. For (d), a product of
singletons has its unique coordinatewise member. Conversely, if the product
has one member, no factor is empty, and replacing one coordinate of that
member by any other value shows that each factor is a singleton. For (e),
two product tuples differ in some coordinate and either tuple supplies a
member of the complementary product. Conversely, combine two distinct
values in \(F_i\) with one complementary tuple.
\end{proof}

\begin{remark}[Exact scope of the fixed-point phrase]
The unconditional obstruction is componentwise: one empty \(F_i\) kills the
global product. The converse claim that no selector implies an empty fixed
locus is valid for finite \(I\), under a supplied product element, or under
a declared suitable choice principle; it is not asserted without that
qualification. A multiple fixed locus alone also does not imply global
multiplicity when another component has an empty fixed locus.
\end{remark}

\section{Relation-image theorem}

Suppose the unchanged core also carries a functor \(E_A:\G\to\mathbf{Set}\)
and a natural transformation \(K:S_\chi\Rightarrow E_A\). Define
\[
  K_*:\Sel(\G)\longrightarrow\Nat(1_{\G},E_A),
  \qquad K_*(\sigma):=K\circ\sigma,
\]
and \(R_i:=K_{r_i}(F_i)\).

\begin{theorem}[Relation-image uniqueness]
\label{thm:kimage}
The map \(K_*\) is well-defined. If \(\Sel(\G)\neq\varnothing\), then
\(\lvert\Img K_*\rvert=1\) if and only if every \(R_i\) is a singleton.
In particular, if some \(F_i\) contains two points with distinct
\(K_{r_i}\)-images, then two global selectors induce different natural
relation sections; a sufficient condition is that \(K_{r_i}\) is injective
on a multiple fixed locus. If every \(R_i\) is a singleton, selector
multiplicity is relation-level choice irrelevance rather than relation
nonuniqueness.
\end{theorem}

\begin{proof}
Naturality of \(K\) and \(\sigma\) makes their composite natural. Fix one
global selector, hence one product tuple, as a baseline. If every \(R_i\) is
a singleton, any two selector tuples have identical \(K\)-images at all
representatives and therefore at all objects. Conversely, if
\(K_{r_i}(x_i)\neq K_{r_i}(x_i')\) for two points of one \(F_i\), replace
only the \(i\)-coordinate of the baseline tuple. Theorem
\ref{thm:classification} reconstructs two selectors whose induced relation
sections differ at \(r_i\).
\end{proof}

The theorem does not assert
\(\Img K_*=\prod_iR_i\) for an arbitrary infinite family without the needed
simultaneous fibre choices. Nor does it assert that every natural relation
section lies in \(\Img K_*\). If \(\Sel(\G)=\varnothing\), its image is
empty, not a uniquely selected relation.

\section{Noninvertible-arrow coherence}

\begin{theorem}[Full-category equalizer guard]
\label{thm:equalizer}
Assume, only conditionally, that \(S_\chi\) extends to a functor
\(\overline S_\chi:\C\to\mathbf{Set}\) whose restriction to the wide core
is \(S_\chi\). Then restriction is injective and identifies
\(\Nat(1_{\C},\overline S_\chi)\) with the subset of
\(\prod_iF_i\) satisfying
\[
  \overline S_\chi(f)
  \bigl(S_\chi(\tau_X)x_{[X]}\bigr)
  =S_\chi(\tau_Y)x_{[Y]}
\]
for every arrow \(f:X\to Y\) of \(\C\), including every noninvertible
arrow. Thus core fixedness is necessary but need not be sufficient for
full-category naturality.
\end{theorem}

\begin{proof}
The core is wide, so restriction retains every object component and is
injective. Theorem \ref{thm:classification} expresses each restricted
selector through its representative tuple. Naturality for the remaining
arrows is exactly the displayed family of equations, proving both inclusion
and sufficiency.
\end{proof}

No \(K\)-claim is extended beyond the core without an explicitly declared
noninvertible relation-transport law. P2-T01 did not adopt a neutral such law
for arbitrary noninvertible maps.

\section{Group actions and anonymous finite structures}

\begin{corollary}[One-object group-action form]
For a group \(G\), its one-object groupoid \(BG\), and a \(G\)-set \(D\),
\[
  \Nat(1_{BG},D)\cong D^G.
\]
Thus the categorical statement and the invariant-point statement are the
same theorem in the one-object case.
\end{corollary}

\begin{corollary}[Transitive action]
\label{cor:transitive}
If a nonempty \(G\)-set \(D\) is transitive and \(|D|>1\), then
\(D^G=\varnothing\), so no deterministic invariant point selector exists.
A transitive singleton has exactly one selector.
\end{corollary}

\begin{proof}
A fixed point has singleton orbit. Transitivity makes its orbit all of
\(D\), contradicting \(|D|>1\).
\end{proof}

\begin{corollary}[Anonymous finite-structure form]
Let \(X\) be a fixed finite source structure with no distinguished
selector-relevant mark, and let its declared structural automorphism group
act transitively on an eligible finite candidate set \(D_X\) with
\(|D_X|>1\). No deterministic isomorphism-invariant point selector exists
on that component.
\end{corollary}

This last statement is only about the declared structural action. It is not
a distributed-computing leader-election theorem, a randomized-selection
result, a dynamical no-go, an empirical claim, or a physical-gauge claim.

\section{Bounded subsumption of prior obstructions}

The following are direct finite instances of Corollary
\ref{cor:transitive}; their full provenance-rich integration remains owned by
P2-T04.
\begin{enumerate}
  \item \textbf{Orientation/line control.} \(\mathrm{GL}(2,2)\) has six
        elements and acts transitively on the three one-dimensional subspaces
        of \(\mathbb F_2^2\). The fixed locus is empty, so no invariant line
        selector exists on that declared component.
  \item \textbf{Directed four-cycle root control.} The four rotations of
        \(C_4\) act transitively on four root choices. The fixed locus is
        empty, so no natural root selector exists on the unmarked cycle.
  \item \textbf{Untagged finite-toy sign/token control.} Sign reversal on
        \(\{+1,-1\}\), and separately a nontrivial token swap on the two
        corresponding token assignments, fixes the untagged source record
        while exchanging every eligible response choice. Each fixed locus is
        empty, so a no-new-source-data equivariant totalization is obstructed.
\end{enumerate}
A trivial two-choice action supplies the complementary nonuniqueness control:
the fixed locus has two elements; an injective \(K\) yields two relation
sections, whereas a constant \(K\) collapses them to one.

\section{Exact limitations and Distance-to-GR status}

The theorem does not prohibit prospectively declared source extensions,
dynamical state selection, randomized selection or invariant probability
measures, or quotient-level observables. A target-neutral extension may
change an automorphism group or fixed locus, but it is new source-extension
data and does not refute the theorem on the original domain. An extension
that directly records the desired answer merely relocates the selection
burden.

No current-ontology derivation or adoption, physical gauge identification,
family reopening or rejection, general \(EqSrc\) discharge, \(M_{\rm src}\)
or \(g_{\rm eff}\) construction, matter coupling, Einstein-equation
derivation, exact-GR recovery, benchmark promotion, Gate Chair verdict,
publication authority, or completed derivation follows. P2-T06 owns the
independent smuggling audit and P2-T07 owns the independent stress test.
The target milestone remains \texttt{source\_equivalence\_eqsrc}; the result
is a precise theorem-level obstruction with no Distance-to-GR ledger delta.

\section*{Reference}

Eilenberg, S., \& Mac Lane, S. (1945). General theory of natural
equivalences. \emph{Transactions of the American Mathematical Society, 58},
231--294. \url{https://doi.org/10.1090/S0002-9947-1945-0013131-6}

\end{document}
